% --- Archon physics formalization source begin ---
% archon:physics
% archon:covers PhyXMiniProblems/problem_phyx_mini_0999.lean
% archon:source-report reports/phyx_mini/problem_phyx_mini_0999.source.json

\chapter{Physics problem phyx\_mini\_0999}
\label{ch:PhyXMiniProblems_problem_phyx_mini_0999}

\paragraph{Problem source.}
\#\# Physical scenario

An imaginary cubical surface of side \textbackslash{}( L \textbackslash{}) is in a region of uniform electric field \textbackslash{}( E \textbackslash{}).

\#\# Question

Find the electric flux through each face of the cube and the total flux through the cube when the cube is turned by an angle \textbackslash{}( 	heta \textbackslash{}) about a vertical axis.

\#\# Answer choices

- A: A: EA \textbackslash{}cos \textbackslash{}theta

- B: B: 0

- C: C: 6EA \textbackslash{}cos \textbackslash{}theta

- D: D: EA(1 - \textbackslash{}cos \textbackslash{}theta)

\#\# Figure caption (auxiliary; use the image as primary evidence)

The image depicts a cube with several arrows and labels. Here's a detailed breakdown:

1. **Cube Orientation:**
   - The cube is shown in three dimensions with dashed lines indicating edges not visible directly.

2. **Surface Normals:**
   - Six arrows labeled \textbackslash{}( \textbackslash{}hat\{n\}\_1 \textbackslash{}) through \textbackslash{}( \textbackslash{}hat\{n\}\_6 \textbackslash{}) represent normal vectors perpendicular to each surface of the cube:
     - \textbackslash{}( \textbackslash{}hat\{n\}\_1 \textbackslash{}), \textbackslash{}( \textbackslash{}hat\{n\}\_4 \textbackslash{}), \textbackslash{}( \textbackslash{}hat\{n\}\_6 \textbackslash{}) point downward or to the side.
     - \textbackslash{}( \textbackslash{}hat\{n\}\_2 \textbackslash{}), \textbackslash{}( \textbackslash{}hat\{n\}\_3 \textbackslash{}), \textbackslash{}( \textbackslash{}hat\{n\}\_5 \textbackslash{}) point upward or outward, indicating the direction normal to the respective faces.

3. **Electric Field Vector (\textbackslash{}( \textbackslash{}vec\{E\} \textbackslash{})):**
   - Pink arrows indicate the direction of an electric field penetrating the cube from one side to the other.

4. **Angles:**
   - The angle \textbackslash{}( \textbackslash{}theta \textbackslash{}) is formed between \textbackslash{}( \textbackslash{}hat\{n\}\_1 \textbackslash{}) and the direction of the electric field (\textbackslash{}( \textbackslash{}vec\{E\} \textbackslash{})).
   - There is a complementary angle labeled \textbackslash{}( 90\textasciicircum{}\textbackslash{}circ - \textbackslash{}theta \textbackslash{}).

5. **Text Elements:**
   - Labels such as \textbackslash{}( \textbackslash{}hat\{n\}\_1, \textbackslash{}hat\{n\}\_2, \textbackslash{}ldots \textbackslash{}hat\{n\}\_6 \textbackslash{}) accompany each normal vector.
   - \textbackslash{}( \textbackslash{}vec\{E\} \textbackslash{}) is labeled next to the set of pink arrows, indicating the presence of an electric field.
   - The angles \textbackslash{}( \textbackslash{}theta \textbackslash{}) and \textbackslash{}( 90\textasciicircum{}\textbackslash{}circ - \textbackslash{}theta \textbackslash{}) are labeled near their respective vector placements.

The image seems to illustrate concepts related to electric flux through a surface, possibly from Gauss's Law in electromagnetism.

\#\# Dataset metadata

Electrostatics; Spatial Relation Reasoning, Physical Model Grounding Reasoning

\paragraph{Recorded answer/context.}
B: B: 0

\paragraph{Figure/image path.}
/root/proposal\_for\_physic/science-mango/phyx\_mini\_run/phyx\_data/test\_image/999.png

\paragraph{Formalization target.}
create a compiling Lean file with sorry bodies at `PhyXMiniProblems/problem\_phyx\_mini\_0999.lean`. The Lean declarations must preserve the physical quantities, dimensions or dimensional roles, figure labels, governing-law hypotheses, and final relation expressed by this problem.
Use Mathlib/Physlib names found through LeanExplore where available. If a physics API is missing, introduce faithful local abstractions rather than scalar placeholder aliases.

\begin{theorem}[Physics formalization target]
\label{thm:physics:phyx_mini_0999:target}
\lean{PhyXMiniProblems.ProblemPhyXMini0999.problem_phyx_mini_0999}
\uses{def:physics:phyx-mini-0999:phyxminiproblems-problemphyxmini0999-lengthinmeters, def:physics:phyx-mini-0999:phyxminiproblems-problemphyxmini0999-fieldmagnitudeinnewtonspercoulomb, def:physics:phyx-mini-0999:phyxminiproblems-problemphyxmini0999-electricfluxinnewtonsquaremeterspercoulomb, def:physics:phyx-mini-0999:phyxminiproblems-problemphyxmini0999-rotatedcubeinuniformfieldsetup, def:physics:phyx-mini-0999:phyxminiproblems-problemphyxmini0999-hasphysicalcubeandfieldparameters, def:physics:phyx-mini-0999:phyxminiproblems-problemphyxmini0999-matchesprimaryrotatedcubefigure, def:physics:phyx-mini-0999:phyxminiproblems-problemphyxmini0999-satisfiesrotatedcubegeometry, def:physics:phyx-mini-0999:phyxminiproblems-problemphyxmini0999-isuniformelectricfieldoncube, def:physics:phyx-mini-0999:phyxminiproblems-problemphyxmini0999-satisfieselectricfluxlaw, def:physics:phyx-mini-0999:phyxminiproblems-problemphyxmini0999-recordeddatasetanswer}
The assigned autoformalize agent should translate this physics problem into Lean declarations in the covered file, with theorem and lemma proof bodies written as `by sorry`.
\end{theorem}
\begin{proof}
This is an autoformalization task, not a proof task. Produce statements that can later be proved by the physics prover without weakening the physical model.
\end{proof}

% --- Archon declaration topology begin ---
\section{Declaration topology}
\begin{definition}[electric Field Strength Dimension]
  \label{def:physics:phyx-mini-0999:phyxminiproblems-problemphyxmini0999-electricfieldstrengthdimension}
  \lean{PhyXMiniProblems.ProblemPhyXMini0999.electricFieldStrengthDimension}
  The physical dimension M L T\textasciicircum{}-2 C\textasciicircum{}-1 of electric-field strength
\end{definition}
\begin{proof}
  This is the declared model definition of electric Field Strength Dimension.
\end{proof}
\begin{definition}[electric Flux Dimension]
  \label{def:physics:phyx-mini-0999:phyxminiproblems-problemphyxmini0999-electricfluxdimension}
  \lean{PhyXMiniProblems.ProblemPhyXMini0999.electricFluxDimension}
  \uses{def:physics:phyx-mini-0999:phyxminiproblems-problemphyxmini0999-electricfieldstrengthdimension}
  The dimension M L\textasciicircum{}3 T\textasciicircum{}-2 C\textasciicircum{}-1 of electric flux E * area
\end{definition}
\begin{proof}
  This is the declared model definition of electric Flux Dimension.
\end{proof}
\begin{definition}[Length Quantity]
  \label{def:physics:phyx-mini-0999:phyxminiproblems-problemphyxmini0999-lengthquantity}
  \lean{PhyXMiniProblems.ProblemPhyXMini0999.LengthQuantity}
  A nonnegative, unit-independent physical length
\end{definition}
\begin{proof}
  This is the declared model definition of Length Quantity.
\end{proof}
\begin{definition}[Area Quantity]
  \label{def:physics:phyx-mini-0999:phyxminiproblems-problemphyxmini0999-areaquantity}
  \lean{PhyXMiniProblems.ProblemPhyXMini0999.AreaQuantity}
  A nonnegative, unit-independent physical area
\end{definition}
\begin{proof}
  This is the declared model definition of Area Quantity.
\end{proof}
\begin{definition}[Electric Field Magnitude Quantity]
  \label{def:physics:phyx-mini-0999:phyxminiproblems-problemphyxmini0999-electricfieldmagnitudequantity}
  \lean{PhyXMiniProblems.ProblemPhyXMini0999.ElectricFieldMagnitudeQuantity}
  \uses{def:physics:phyx-mini-0999:phyxminiproblems-problemphyxmini0999-electricfieldstrengthdimension}
  A nonnegative, unit-independent electric-field magnitude
\end{definition}
\begin{proof}
  This is the declared model definition of Electric Field Magnitude Quantity.
\end{proof}
\begin{definition}[Signed Electric Flux Quantity]
  \label{def:physics:phyx-mini-0999:phyxminiproblems-problemphyxmini0999-signedelectricfluxquantity}
  \lean{PhyXMiniProblems.ProblemPhyXMini0999.SignedElectricFluxQuantity}
  \uses{def:physics:phyx-mini-0999:phyxminiproblems-problemphyxmini0999-electricfluxdimension}
  A signed outward electric flux through a face or closed surface
\end{definition}
\begin{proof}
  This is the declared model definition of Signed Electric Flux Quantity.
\end{proof}
\begin{definition}[length In Meters]
  \label{def:physics:phyx-mini-0999:phyxminiproblems-problemphyxmini0999-lengthinmeters}
  \lean{PhyXMiniProblems.ProblemPhyXMini0999.lengthInMeters}
  \uses{def:physics:phyx-mini-0999:phyxminiproblems-problemphyxmini0999-lengthquantity}
  Coherent-SI readout of a physical length, in metres
\end{definition}
\begin{proof}
  This is the declared model definition of length In Meters.
\end{proof}
\begin{definition}[area In Square Meters]
  \label{def:physics:phyx-mini-0999:phyxminiproblems-problemphyxmini0999-areainsquaremeters}
  \lean{PhyXMiniProblems.ProblemPhyXMini0999.areaInSquareMeters}
  \uses{def:physics:phyx-mini-0999:phyxminiproblems-problemphyxmini0999-areaquantity}
  Coherent-SI readout of a physical area, in square metres
\end{definition}
\begin{proof}
  This is the declared model definition of area In Square Meters.
\end{proof}
\begin{definition}[field Magnitude In Newtons Per Coulomb]
  \label{def:physics:phyx-mini-0999:phyxminiproblems-problemphyxmini0999-fieldmagnitudeinnewtonspercoulomb}
  \lean{PhyXMiniProblems.ProblemPhyXMini0999.fieldMagnitudeInNewtonsPerCoulomb}
  \uses{def:physics:phyx-mini-0999:phyxminiproblems-problemphyxmini0999-electricfieldmagnitudequantity}
  Coherent-SI readout of field magnitude, in newtons per coulomb
\end{definition}
\begin{proof}
  This is the declared model definition of field Magnitude In Newtons Per Coulomb.
\end{proof}
\begin{definition}[electric Flux In Newton Square Meters Per Coulomb]
  \label{def:physics:phyx-mini-0999:phyxminiproblems-problemphyxmini0999-electricfluxinnewtonsquaremeterspercoulomb}
  \lean{PhyXMiniProblems.ProblemPhyXMini0999.electricFluxInNewtonSquareMetersPerCoulomb}
  \uses{def:physics:phyx-mini-0999:phyxminiproblems-problemphyxmini0999-signedelectricfluxquantity}
  Coherent-SI readout of signed flux, in newton-square-metres per coulomb
\end{definition}
\begin{proof}
  This is the declared model definition of electric Flux In Newton Square Meters Per Coulomb.
\end{proof}
\begin{definition}[Cube Face]
  \label{def:physics:phyx-mini-0999:phyxminiproblems-problemphyxmini0999-cubeface}
  \lean{PhyXMiniProblems.ProblemPhyXMini0999.CubeFace}
  The six outward-normal labels printed in the primary raster
\end{definition}
\begin{proof}
  This is the declared model definition of Cube Face.
\end{proof}
\begin{definition}[Rotated Cube Figure]
  \label{def:physics:phyx-mini-0999:phyxminiproblems-problemphyxmini0999-rotatedcubefigure}
  \lean{PhyXMiniProblems.ProblemPhyXMini0999.RotatedCubeFigure}
  \uses{def:physics:phyx-mini-0999:phyxminiproblems-problemphyxmini0999-cubeface}
  Typed content of the supplied raster. Direction vectors are dimensionless; physical field magnitude is stored independently in the setup
\end{definition}
\begin{proof}
  This is the declared model definition of Rotated Cube Figure.
\end{proof}
\begin{definition}[Rotated Cube In Uniform Field Setup]
  \label{def:physics:phyx-mini-0999:phyxminiproblems-problemphyxmini0999-rotatedcubeinuniformfieldsetup}
  \lean{PhyXMiniProblems.ProblemPhyXMini0999.RotatedCubeInUniformFieldSetup}
  \uses{def:physics:phyx-mini-0999:phyxminiproblems-problemphyxmini0999-lengthquantity, def:physics:phyx-mini-0999:phyxminiproblems-problemphyxmini0999-areaquantity, def:physics:phyx-mini-0999:phyxminiproblems-problemphyxmini0999-electricfieldmagnitudequantity, def:physics:phyx-mini-0999:phyxminiproblems-problemphyxmini0999-signedelectricfluxquantity, def:physics:phyx-mini-0999:phyxminiproblems-problemphyxmini0999-cubeface, def:physics:phyx-mini-0999:phyxminiproblems-problemphyxmini0999-rotatedcubefigure}
  Independent physical data for the rotated cube. The face and total fluxes are unconstrained physical quantities here; their values are related to the field only by SatisfiesElectricFluxLaw below
\end{definition}
\begin{proof}
  This is the declared model definition of Rotated Cube In Uniform Field Setup.
\end{proof}
\begin{definition}[Has Physical Cube And Field Parameters]
  \label{def:physics:phyx-mini-0999:phyxminiproblems-problemphyxmini0999-hasphysicalcubeandfieldparameters}
  \lean{PhyXMiniProblems.ProblemPhyXMini0999.HasPhysicalCubeAndFieldParameters}
  \uses{def:physics:phyx-mini-0999:phyxminiproblems-problemphyxmini0999-lengthinmeters, def:physics:phyx-mini-0999:phyxminiproblems-problemphyxmini0999-areainsquaremeters, def:physics:phyx-mini-0999:phyxminiproblems-problemphyxmini0999-fieldmagnitudeinnewtonspercoulomb, def:physics:phyx-mini-0999:phyxminiproblems-problemphyxmini0999-rotatedcubeinuniformfieldsetup}
  Positivity conditions for the physical parameters named in the problem
\end{definition}
\begin{proof}
  This is the declared model definition of Has Physical Cube And Field Parameters.
\end{proof}
\begin{definition}[Matches Primary Rotated Cube Figure]
  \label{def:physics:phyx-mini-0999:phyxminiproblems-problemphyxmini0999-matchesprimaryrotatedcubefigure}
  \lean{PhyXMiniProblems.ProblemPhyXMini0999.MatchesPrimaryRotatedCubeFigure}
  \uses{def:physics:phyx-mini-0999:phyxminiproblems-problemphyxmini0999-rotatedcubeinuniformfieldsetup}
  Primary-image evidence. It records all six normal labels, the family of parallel field arrows, and the two printed angle annotations. It contains no flux value
\end{definition}
\begin{proof}
  This is the declared model definition of Matches Primary Rotated Cube Figure.
\end{proof}
\begin{definition}[Satisfies Rotated Cube Geometry]
  \label{def:physics:phyx-mini-0999:phyxminiproblems-problemphyxmini0999-satisfiesrotatedcubegeometry}
  \lean{PhyXMiniProblems.ProblemPhyXMini0999.SatisfiesRotatedCubeGeometry}
  \uses{def:physics:phyx-mini-0999:phyxminiproblems-problemphyxmini0999-lengthinmeters, def:physics:phyx-mini-0999:phyxminiproblems-problemphyxmini0999-areainsquaremeters, def:physics:phyx-mini-0999:phyxminiproblems-problemphyxmini0999-rotatedcubeinuniformfieldsetup}
  Geometry of the rotated cube. The component equations formalize the two angles read from the raster: n1 makes angle theta with the direction opposite the field, while n4 makes angle pi/2 - theta with the field. They are geometric facts about unit directions, not assumptions about flux
\end{definition}
\begin{proof}
  This is the declared model definition of Satisfies Rotated Cube Geometry.
\end{proof}
\begin{definition}[Is Uniform Electric Field On Cube]
  \label{def:physics:phyx-mini-0999:phyxminiproblems-problemphyxmini0999-isuniformelectricfieldoncube}
  \lean{PhyXMiniProblems.ProblemPhyXMini0999.IsUniformElectricFieldOnCube}
  \uses{def:physics:phyx-mini-0999:phyxminiproblems-problemphyxmini0999-fieldmagnitudeinnewtonspercoulomb, def:physics:phyx-mini-0999:phyxminiproblems-problemphyxmini0999-rotatedcubeinuniformfieldsetup}
  The stated electric field is uniform throughout the region containing the cube. The Physlib field vector is calibrated to the independent dimensionful magnitude and unit direction in coherent SI components
\end{definition}
\begin{proof}
  This is the declared model definition of Is Uniform Electric Field On Cube.
\end{proof}
\begin{definition}[Satisfies Electric Flux Law]
  \label{def:physics:phyx-mini-0999:phyxminiproblems-problemphyxmini0999-satisfieselectricfluxlaw}
  \lean{PhyXMiniProblems.ProblemPhyXMini0999.SatisfiesElectricFluxLaw}
  \uses{def:physics:phyx-mini-0999:phyxminiproblems-problemphyxmini0999-areainsquaremeters, def:physics:phyx-mini-0999:phyxminiproblems-problemphyxmini0999-electricfluxinnewtonsquaremeterspercoulomb, def:physics:phyx-mini-0999:phyxminiproblems-problemphyxmini0999-cubeface, def:physics:phyx-mini-0999:phyxminiproblems-problemphyxmini0999-rotatedcubeinuniformfieldsetup}
  The planar-surface electric-flux law and finite additivity over the closed cube. Neither field states any requested trigonometric face formula or the zero-total-flux conclusion
\end{definition}
\begin{proof}
  This is the declared model definition of Satisfies Electric Flux Law.
\end{proof}
\begin{definition}[Answer Choice]
  \label{def:physics:phyx-mini-0999:phyxminiproblems-problemphyxmini0999-answerchoice}
  \lean{PhyXMiniProblems.ProblemPhyXMini0999.AnswerChoice}
  Labels of the four electric-flux choices printed in the source
\end{definition}
\begin{proof}
  This is the declared model definition of Answer Choice.
\end{proof}
\begin{definition}[displayed Flux In Newton Square Meters Per Coulomb]
  \label{def:physics:phyx-mini-0999:phyxminiproblems-problemphyxmini0999-answerchoice-displayedfluxinnewtonsquaremeterspercoulomb}
  \lean{PhyXMiniProblems.ProblemPhyXMini0999.AnswerChoice.displayedFluxInNewtonSquareMetersPerCoulomb}
  \uses{def:physics:phyx-mini-0999:phyxminiproblems-problemphyxmini0999-lengthinmeters, def:physics:phyx-mini-0999:phyxminiproblems-problemphyxmini0999-fieldmagnitudeinnewtonspercoulomb, def:physics:phyx-mini-0999:phyxminiproblems-problemphyxmini0999-rotatedcubeinuniformfieldsetup, def:physics:phyx-mini-0999:phyxminiproblems-problemphyxmini0999-answerchoice}
  Flux readout printed beside each answer label, with the source's A interpreted as the square face area L\textasciicircum{}2
\end{definition}
\begin{proof}
  This is the declared model definition of displayed Flux In Newton Square Meters Per Coulomb.
\end{proof}
\begin{definition}[recorded Dataset Answer]
  \label{def:physics:phyx-mini-0999:phyxminiproblems-problemphyxmini0999-recordeddatasetanswer}
  \lean{PhyXMiniProblems.ProblemPhyXMini0999.recordedDatasetAnswer}
  \uses{def:physics:phyx-mini-0999:phyxminiproblems-problemphyxmini0999-answerchoice}
  The answer label recorded in the dataset
\end{definition}
\begin{proof}
  This is the declared model definition of recorded Dataset Answer.
\end{proof}
\begin{lemma}[flux through each face]
  \label{lem:physics:phyx-mini-0999:phyxminiproblems-problemphyxmini0999-flux-through-each-face}
  \lean{PhyXMiniProblems.ProblemPhyXMini0999.flux_through_each_face}
  \uses{def:physics:phyx-mini-0999:phyxminiproblems-problemphyxmini0999-lengthinmeters, def:physics:phyx-mini-0999:phyxminiproblems-problemphyxmini0999-fieldmagnitudeinnewtonspercoulomb, def:physics:phyx-mini-0999:phyxminiproblems-problemphyxmini0999-electricfluxinnewtonsquaremeterspercoulomb, def:physics:phyx-mini-0999:phyxminiproblems-problemphyxmini0999-rotatedcubeinuniformfieldsetup, def:physics:phyx-mini-0999:phyxminiproblems-problemphyxmini0999-satisfiesrotatedcubegeometry, def:physics:phyx-mini-0999:phyxminiproblems-problemphyxmini0999-isuniformelectricfieldoncube, def:physics:phyx-mini-0999:phyxminiproblems-problemphyxmini0999-satisfieselectricfluxlaw}
  The six requested signed outward face fluxes. Faces n1 and n2 cancel, as do n3 and n4; the vertical faces n5 and n6 carry no flux
\end{lemma}
\begin{proof}
  The conclusion follows from the governing relations and figure readouts named in the statement of flux through each face.
\end{proof}
\begin{lemma}[total outward flux is zero]
  \label{lem:physics:phyx-mini-0999:phyxminiproblems-problemphyxmini0999-total-outward-flux-is-zero}
  \lean{PhyXMiniProblems.ProblemPhyXMini0999.total_outward_flux_is_zero}
  \uses{def:physics:phyx-mini-0999:phyxminiproblems-problemphyxmini0999-electricfluxinnewtonsquaremeterspercoulomb, def:physics:phyx-mini-0999:phyxminiproblems-problemphyxmini0999-rotatedcubeinuniformfieldsetup, def:physics:phyx-mini-0999:phyxminiproblems-problemphyxmini0999-satisfiesrotatedcubegeometry, def:physics:phyx-mini-0999:phyxminiproblems-problemphyxmini0999-isuniformelectricfieldoncube, def:physics:phyx-mini-0999:phyxminiproblems-problemphyxmini0999-satisfieselectricfluxlaw}
  Finite additivity and opposite-face cancellation make the total outward flux zero for every turn angle
\end{lemma}
\begin{proof}
  The conclusion follows from the governing relations and figure readouts named in the statement of total outward flux is zero.
\end{proof}
% --- Archon declaration topology end ---
% --- Archon physics formalization source end ---
